% !TeX root = ../main.tex
% ============================================================
%  sections/01_introduction.tex
% ============================================================

\section{Introduction}
\label{sec:introduction}

Climate change and the regulatory pressure for sustainable industrial
practices have prompted governments worldwide to deploy market-based
instruments for controlling greenhouse-gas (GHG) emissions. Among
these, \emph{carbon taxation} is widely regarded as one of the most
transparent and economically efficient policy tools: by attaching a
monetary price to each unit of CO\textsubscript{2}-equivalent
emitted, it internalises the negative externality of pollution and
creates incentives for supply chains to redesign networks, adopt
cleaner production technologies, and invest in reverse logistics
\citep{nordhaus2007, pigou1920}. Evidence from an implemented carbon
tax shows that a well-calibrated rate can produce a substantial,
empirically verified drop in emissions \citep{andersson2019}, although
the effectiveness of any given rate is mediated by the structure of
the affected supply chains.

Supply chain network design (\SCND) is one of the most consequential
strategic decisions an organisation makes. Facility location, capacity
sizing, technology installation, and transportation contracts
typically lock in operating costs and emissions over horizons of
several decades \citep{melo2009}. When a carbon tax is in force,
these long-lived decisions must simultaneously balance cost
efficiency, regulatory compliance, and environmental performance.
The rise of \emph{closed-loop supply chains} (\CLSC), which integrate
forward distribution with the recovery, remanufacturing, and reuse of
end-of-life items, further expands the design space and enlarges the
opportunities to abate emissions through circular operations
\citep{fleischmann2001, govindan2015}.

Crucially, the interaction between the regulator and the supply chain
operator is \emph{inherently hierarchical}: the regulator commits
first to a tax level, after which the operator optimises its network
and operations. Conventional single-level \SCND models, including
emission-aware variants, take the carbon tax as a fixed parameter and
therefore cannot represent the regulator's strategic problem of
\emph{choosing} a tax level so as to induce a desirable response
from rational, cost-minimising supply chain operators. As a
consequence, they may
recommend tax rates that are either ineffective (failing to trigger
network reconfiguration) or excessive (imposing economic costs out
of proportion to the abatement they achieve).

\subsection{Motivation}
\label{sec:intro:motivation}

Bilevel optimisation provides a natural framework for modelling
Stackelberg leader--follower games \citep{colson2007, dempe2002}. In
the present setting:
\begin{itemize}
  \item the \textbf{leader} (regulator) selects the carbon tax, with
        the goal of minimising the supply chain's total GHG emissions
        and, among all tax levels that achieve that minimum, levying
        the lowest one;
  \item the \textbf{follower} (supply chain operator) solves a
        multi-echelon, multi-period \CLSC design and operations
        problem parameterised by the carbon tax, minimising
        total operating costs inclusive of tax payments.
\end{itemize}
The resulting \BLMIP is NP-hard even when both levels are linear
programs \citep{jeroslow1985}. The presence of binary facility,
technology, and capacity decisions in the follower compounds this
difficulty and rules out direct application of techniques designed
for continuous bilevel programs. Specialised solution methods are
therefore required.

\subsection{Contributions}
\label{sec:intro:contributions}

This paper makes three interrelated contributions, each positioned
against the closest related work in \Cref{sec:lit:gap}:
\begin{enumerate}
  \item We formulate a tactical \BLMIP in which a regulator
        chooses the carbon tax from a finite grid --- minimising
        induced emissions, then the tax itself --- anticipating the
        optimal response of a multi-echelon, multi-period closed-loop
        supply chain operator whose facility, technology, and
        capacity decisions remain integer.
  \item We develop a \DBP that exploits the leader's
        finite tax grid directly, combining valid lower bounds from a
        parametric \LP relaxation of the follower with valid upper
        bounds from exact follower responses at sampled tax levels,
        with proved finite-time convergence to a global
        $\epsilon$-optimal solution.
  \item Runs
        on generated benchmark instances show that the optimal tax is
        strongly instance-specific, that the emissions response
        saturates beyond a threshold --- an over-taxation pathology
        for a purely emissions-minimising regulator --- and that
        cleaner production/remanufacturing technology, not
        transport-mode choice, is the binding abatement lever.
\end{enumerate}

\subsection{Paper organisation}
\label{sec:intro:organisation}

The remainder of the paper is structured as follows.
\Cref{sec:literature} reviews the related literature on bilevel
optimisation, sustainable \SCND, and carbon pricing in operations
management, and positions our contribution.
\Cref{sec:problem} describes the physical network and the
hierarchical decision structure.
\Cref{sec:model} states the bilevel mixed-integer mathematical
formulation.
\Cref{sec:methodology} presents the \DBP and establishes its
convergence.
\Cref{sec:experiments} reports computational experiments.
\Cref{sec:discussion} discusses the managerial and policy
implications, robustness of the findings, and limitations.
\Cref{sec:conclusion} concludes and outlines directions for
further research.
